% notation.tex — THE symbol registry. Defined ONCE, used everywhere (book skill).
% Extend this file; never redefine a symbol locally in a chapter. Seeded here with
% what Chapter 1 (Philosophy) needs; later chapters extend it (architecture, then
% the mathematics chapter, which owns the coboundary/grounding symbols).

% --- Proper names (used throughout). --------------------------------------------
% The framework vs. the live system it runs (memory `ouroboros-naming`; the live
% system is defined in \artifact{dn-ouroboros-principal} §1).
\newcommand{\mindpalace}{\textsc{mind-palace}}      % the framework: repo + machinery
\newcommand{\Ouroboros}{\textsc{Ouroboros}}          % the live system: daemon + corpus
\newcommand{\Constitution}{\textsc{Constitution}}    % the fixed-point kernel

% --- Artifact-chain node names (shared by prose and the Fig. 1.1 TikZ source). ---
% Grounded in \artifact{dn-agent-workflow} §3 (the state-machine table).
\newcommand{\stbrainstorm}{\textsf{brainstorm}}
\newcommand{\stdesign}{\textsf{design note}}
\newcommand{\stplan}{\textsf{build plan}}
\newcommand{\stjournal}{\textsf{journal + findings}}
\newcommand{\streflect}{\textsf{reflection}}

% =============================================================================
% ADDED BY CHAPTER 2 (Architecture), bp-104. Every symbol below is grounded in a
% ratified source, named in the comment. Extend this file; never fork it.
% =============================================================================

% --- The trust zones (\coderef{docs/BUILD-SPEC.md}{\gitref} §6). -----------------
\newcommand{\zoneA}{\textsf{Zone\,A}}                % sealed core, no network
\newcommand{\zoneB}{\textsf{Zone\,B}}                % networked edge
\newcommand{\zoneC}{\textsf{Zone\,C}}                % cloud (AWS)

% --- The four OS plane principals (\artifact{dn-plane-principals} §3.1). ---------
\newcommand{\phuman}{\texttt{ascalva}}               % the human login
\newcommand{\pwork}{\texttt{ouroboros-work}}         % the workflow plane
\newcommand{\pcore}{\texttt{ouroboros}}              % the core plane
\newcommand{\pedge}{\texttt{ouroboros-edge}}         % the edge plane

% --- The import-closure boundary (\artifact{dn-inner-outer-core} §2.1, §2.3;
%     \coderef{ops/import\_lint.py}{\gitref}). -----------------------------------
\newcommand{\coremods}{\mathcal{C}}                  % the modules under core/
\newcommand{\netmods}{\mathcal{N}}                   % network-capable module roots
\newcommand{\netallow}{\mathcal{L}}                  % the audited loopback allowlist
\newcommand{\basemods}{\mathcal{B}}                  % the admissible base of the ring predicate
\newcommand{\innerring}{\mathcal{I}}                 % the inner ring (the computed fixed point)
\newcommand{\imports}{\operatorname{imp}}            % direct imports of a module
\newcommand{\reach}{\operatorname{Reach}}            % transitive import reachability

% --- The capability-scope algebra (\artifact{dn-capability-scope} §2.1--§2.3). ---
% A scope is s = (Sigma, E, T, A). The symbols are the note's own; do not rename.
\newcommand{\scopeobj}{s}                            % a scope
\newcommand{\stratums}{\Sigma}                       % stratum scope
\newcommand{\edgeclass}{E}                           % edge-class (fiber) scope
\newcommand{\timescope}{T}                           % time scope = (clock, window)
\newcommand{\authority}{A}                           % authority = P x W_Sigma x W_world
\newcommand{\denyset}{\mathfrak{D}}                  % the foundation denylist, as an order ideal
\newcommand{\reffor}{R}                              % the stratum refinement forest
\newcommand{\meetop}{\sqcap}                         % meet: safe composition
\newcommand{\joinop}{\sqcup}                         % join: a widening
\newcommand{\scopeleq}{\sqsubseteq}                  % scope ordering
\newcommand{\clockk}{\kappa}                         % a clock (a monotone coarsening)
\newcommand{\Inv}{\textsf{Inv}}                      % clock-independent result type
\newcommand{\Rate}{\textsf{Rate}}                    % clock-indexed rate result type

% --- The enforcement-tier ladder (\artifact{dn-capability-scope} §2.2): the
%     chapter's grading rubric. Composition takes the MIN tier along the chain. ---
\newcommand{\tstruct}{\textsf{structural}}
\newcommand{\tstatic}{\textsf{static+guard}}
\newcommand{\tconv}{\textsf{convention}}

% --- Reserved for later chapters (NOT used yet; declared here so later scribe
%     plans EXTEND rather than fork the registry). The mathematics chapter is the
%     canonical home of these; their definitions ratify with that chapter's sources.
%     Kept commented until first use to avoid asserting draft formalism in the book.
% \newcommand{\authdist}{\alpha}              % authorship-distance coordinate (draft)
% \newcommand{\footprint}[1]{\hat{\alpha}(#1)}% derived-node authorship footprint (draft)
% \newcommand{\confbound}{c \le \gamma^{d}\, g}% the grounding/confidence envelope
